%==============================================================================
\chapter{Introduction}
\label{ch:intro}
%==============================================================================

\section{Overview}
\label{sec:overview}

Good explanations are often visual and spoken at the same time. A teacher
points at a part of a drawing and says what it does; the pointing and the
talking happen together, so the listener never has to guess which part the
words are about. Most digital learning tools lose this link. A page of text
sits next to a fixed picture, and a video lives somewhere else. The learner
carries the load of connecting them.

This thesis describes \textbf{TutorAI}, a working end-to-end system that
rebuilds this ``point and explain'' experience automatically. The learner
types a single topic --- for example, \emph{Photosynthesis} --- into an
Android application. The system returns a self-contained \emph{lesson} made of
three parts:

\begin{enumerate}
  \item a diagram drawn as Scalable Vector Graphics (SVG) in which every
        meaningful visual element carries a stable, semantic identifier such
        as \texttt{sun}, \texttt{leaf}, or \texttt{chloroplast};
  \item a narration script split into ordered \emph{segments}, where each
        segment consists of a sentence or two of explanation together with the
        list of SVG element identifiers it talks about; and
  \item one synthesized speech clip per segment, produced by a text-to-speech
        (TTS) engine, with a precisely measured duration.
\end{enumerate}

During playback on the device, the current segment's audio plays while the SVG
elements it references are visually highlighted (a glow-and-scale emphasis),
so the learner always sees exactly what is being described. When the segment's
clip finishes, the highlight moves on with the narration. Because every asset
of a lesson is downloaded once and stored on the device, a saved lesson
replays fully offline, and an on-device history lets the learner revisit any
past topic.

The system consists of two cooperating applications. The \emph{backend} is a
stateless Python FastAPI service that hosts an agentic generation pipeline
built with LangChain and LangGraph. The pipeline orchestrates a Gemini large
language model (LLM) for planning, drawing, and narration, and a Gemini speech
model for audio, both hidden behind swappable provider adapters. Generated
lessons are cached under a normalized topic key so that a popular topic is
generated once and afterwards served instantly and cheaply. The \emph{client}
is a Kotlin Android application built with Jetpack Compose, following Clean
Architecture and the Model--View--ViewModel (MVVM) pattern, with Room and a
private file store for offline persistence and ExoPlayer for per-segment audio
playback.

The intellectual core of the work is not any single model but an
\emph{architectural} answer to a reliability problem: an LLM will happily
write narration that points at a diagram part that does not exist, which
silently breaks the highlighting that the whole product depends on. TutorAI
therefore splits generation into a graph of small steps with bounded
critique-and-refine and repair loops, and places a deterministic, non-LLM
validator as a hard gate in front of the cache. A lesson is never stored
unless its SVG is well formed and every identifier referenced by every
segment actually exists in the drawing.

\section{Problem Statement}
\label{sec:problem}

Learners who want a quick, visual explanation of a concept currently have to
stitch together text, static diagrams, and separately produced videos. There
is no lightweight tool where the user types a topic and immediately receives a
\emph{purpose-built diagram that explains itself} --- narrated, with the
relevant parts of the diagram lighting up exactly as they are discussed ---
and that can be revisited later without a network connection.

Building such a tool with generative AI raises four concrete technical
problems that this thesis addresses:

\begin{enumerate}
  \item \textbf{Cross-artifact consistency.} The diagram, the narration, and
        the audio are generated by different (probabilistic) steps, yet the
        product depends on an exact correspondence between narration
        references and diagram element identifiers. A single hallucinated
        identifier breaks synchronized highlighting.
  \item \textbf{Deterministic audio--visual synchronization.} Highlighting
        must stay aligned with speech on a mobile device, offline, without
        streaming timing data or word-level alignment infrastructure.
  \item \textbf{Cost and latency of generation.} Cold generation requires
        several LLM calls plus per-segment speech synthesis. Repeated requests
        for the same topic must not repeat this cost, and unbounded
        self-correction loops must not inflate it.
  \item \textbf{Reproducible engineering.} The system must be testable without
        live AI keys, deployable as a container, and resilient to externally
        controlled details such as hosted model identifiers changing over
        time.
\end{enumerate}

\section{Objectives}
\label{sec:objectives}

The objectives of this project work are:

\begin{enumerate}
  \item To design and implement an \textbf{agentic generation pipeline}
        (LangGraph state graph) that transforms a free-text topic into a
        validated lesson: an SVG diagram with stable semantic identifiers,
        ordered narration segments that reference only existing identifiers,
        and per-segment TTS audio with measured durations.
  \item To design a \textbf{deterministic, non-LLM validation gate} with
        bounded repair and critique-and-refine loops, guaranteeing that no
        lesson whose narration references a non-existent diagram element is
        ever cached or served.
  \item To implement a \textbf{stateless FastAPI backend} with an
        asynchronous job model, a shared topic-keyed lesson cache with request
        coalescing, an asset store, and swappable LLM/TTS provider adapters.
  \item To implement an \textbf{offline-first Android client} (Kotlin,
        Jetpack Compose, Clean Architecture + MVVM) that generates lessons via
        job polling, plays them with segment-synchronized SVG highlighting,
        and persists them (Room + file store) for fully offline replay with
        an on-device history.
  \item To apply sound \textbf{software-engineering practice} throughout:
        documented UML-first design, classical design patterns at the seams,
        hermetic key-free automated tests, containerization, and CI/CD
        pipelines for both backend and client.
  \item To \textbf{verify the realized system} end to end against live
        models and report its behaviour honestly, distinguishing verified
        results from future evaluation.
\end{enumerate}

\section{Motivation}
\label{sec:motivation}

The motivation for this work is threefold.

\textbf{Pedagogical.} Decades of research on multimedia learning show that
people learn more deeply from words and pictures presented together than from
words alone, and that spoken narration synchronized with relevant graphics
reduces cognitive load compared to on-screen text (the modality and temporal
contiguity principles). Bloom's classic finding that one-to-one tutoring
dramatically outperforms group instruction has driven a long search for
scalable tutoring software. A tool that generates a narrated, self-pointing
diagram for \emph{any} typed topic, on demand, brings a small but genuinely
useful slice of that tutoring behaviour to anyone with a phone.

\textbf{Technical.} Large language models have made it possible to generate
explanations, structured data, and even vector graphics from a single prompt,
but they remain probabilistic: outputs may be malformed or internally
inconsistent. The recent literature on LLM agents proposes decomposition,
self-critique, and iterative repair as remedies, yet published systems rarely
need the \emph{hard} cross-artifact guarantee that TutorAI needs (every
narration reference must resolve to a real diagram element). Demonstrating
that a small state graph with bounded loops and a deterministic gate can make
a probabilistic generator reliable enough for a correctness-critical feature
is a useful, transferable engineering result.

\textbf{Practical.} Learners in many settings have intermittent or expensive
connectivity. An offline-first design --- download once, replay forever,
history on the device, no accounts, no personal data on the server --- makes
the tool usable and private in exactly the settings where lightweight learning
aids matter most. A stateless backend with a shared cache also keeps the
operating cost of such a service low: the marginal cost of a popular topic
tends towards zero.

\section{Organization of the Thesis}
\label{sec:organization}

The remainder of this thesis is organized as follows.
Chapter~\ref{ch:literature} surveys the relevant literature across
intelligent tutoring systems, multimedia learning theory, large language
models and agentic orchestration, automatic diagram generation, speech
synthesis, and offline-first mobile architecture, and identifies the research
gap. Chapter~\ref{ch:methodology} presents the methodology: the requirements,
the key design decisions, the agentic generation method with its validation
gate, the synchronization model, and the phased development and testing
process followed. Chapter~\ref{ch:design} details the design and
implementation of the backend and the Android client, including the domain
model, the API contract, the design patterns employed, and the CI/CD
pipelines. Chapter~\ref{ch:results} reports the verified results of the
realized system and discusses its behaviour, trade-offs, and limitations.
Chapter~\ref{ch:conclusion} concludes and outlines the future scope. The
appendices reproduce the API contract, the lesson manifest schema, and the
repository layout.
